\uuid{nT2w}
\titre{Effet du stride sur la taille et la résolution spatiale}
\niveau{L2}
\module{Informatique / Signal}
\chapitre{Réseaux de neurones}
\sousChapitre{Stride et sous-échantillonnage}
\theme{Calcul de taille, résolution, stride inverse}
\auteur{Maxime Nguyen}
\datecreate{2026-03-23}
\organisation{AMSCC}
\difficulte{2}

\contenu{
\texte{
  On dispose d'une image $6 \times 6$ et d'un noyau moyenneur $3 \times 3$
  (tous les coefficients valent $1/9$), avec padding $p = 1$.
}
\begin{enumerate}
\item
  \question{Calculer la taille de sortie pour les strides $s = 1$, $s = 2$ et $s = 3$.}
  \reponse{
    Avec $\text{in} = 6$, $p = 1$, $k = 3$ :
    \begin{itemize}
      \item $s=1$ : $\lfloor(6+2-3)/1\rfloor+1 = 5+1 = 6$.
      \item $s=2$ : $\lfloor(6+2-3)/2\rfloor+1 = \lfloor5/2\rfloor+1 = 2+1 = 3$.
      \item $s=3$ : $\lfloor(6+2-3)/3\rfloor+1 = \lfloor5/3\rfloor+1 = 1+1 = 2$.
    \end{itemize}
  }

\item
  \question{Pour stride $= 2$ : combien de fois le noyau se déplace-t-il en ligne ? En colonne ?}
  \reponse{
    La sortie est $3 \times 3$ : le noyau occupe $3$ positions distinctes en ligne et
    $3$ en colonne. Il effectue donc $3 \times 3 = 9$ placements au total
    (contre $36$ pour stride $= 1$).
  }

\item
  \question{Deux pixels voisins de l'image originale (distance $1$) correspondent-ils toujours à deux pixels voisins dans la sortie avec stride $= 2$ ? Qu'est-ce que cela implique sur la résolution spatiale ?}
  \reponse{
    Non. Avec stride $= 2$, deux positions consécutives du noyau sont séparées de $2$
    pixels dans l'image d'origine. Un pixel sur deux de l'entrée est donc
    \og sauté \fg{}. La résolution spatiale est divisée par $2$ : des structures
    plus fines que $2$ pixels ne peuvent plus être distinguées dans la sortie.
    C'est un phénomène de \textbf{sous-échantillonnage} (aliasing possible si le
    signal n'est pas suffisamment lisse).
  }

\item
  \question{On souhaite obtenir une sortie de taille $3 \times 3$ à partir d'une image $7 \times 7$, avec un noyau $3 \times 3$ et \textbf{sans padding}. Quel stride faut-il utiliser ?}
  \indication{Exprimer $s$ en fonction de la formule et résoudre.}
  \reponse{
    On cherche $s$ tel que $\lfloor(7+0-3)/s\rfloor+1 = 3$, soit
    $\lfloor4/s\rfloor = 2$, donc $2 \le 4/s < 3$, ce qui donne
    $4/3 < s \le 2$. Le seul entier dans cet intervalle est $s = 2$.

    Vérification : $\lfloor4/2\rfloor+1 = 2+1 = 3$. ✓
  }
\end{enumerate}
}
